\section{The work contract}
For exposition, write an agreement as
\[
 A=(B,P,I,C,D,F,T,L),\qquad a=H(\operatorname{canon}(A)).
\]
$I$ commits the input, $C$ identifies the checker and acceptance policy, $D$
defines permitted disclosure, $F$ identifies the funding domain and terms, $T$
sets the deadlines, and $L$ names any parent agreement. Both parties sign the
same canonical agreement. This tuple is a description of the common contract,
not a replacement wire encoding. The artifact companion maps it to the
versioned implementation objects.

An agreement has three bindings. \emph{Authority binding} connects it to the
receiver's own admitted operation and bounded local grant. \emph{Result binding}
connects the exact input, output and acceptance procedure. \emph{Funding binding}
connects the agreement and result commitment to one allocation, beneficiary,
asset and amount in one settlement domain. These bindings prevent a valid
statement about one relationship from being reused as a statement about another.
They do not establish useful work independently of the acceptance procedure.

\subsection{Funding is an allocation}
The payer deposits $q>0$ units into an allocation $f$. The reference contract
derives its allocation identifier from the chain, contract address and complete
terms. The agreement digest, payer, beneficiary, verifier, asset, amount and
four deadlines are immutable. A new payer agreement cannot be funded twice
under the same payer and digest. The receiver checks the selected contract code
and observation source before relying on an allocation.

A provider does not start merely because a payer supplies a signed balance. It
requires an acceptable observation of the actual reserved allocation and a
local admission decision bound to it. The binding does not atomically fuse a
remote blockchain with a local database. Fresh observation, the deadline
margin, retained operation identity and conservative recovery are necessary.
If the chain reorganizes after an observation, the profile's finality assumption
has failed or the observation was premature. A local receipt cannot repair that.

\subsection{Acceptance is an explicit decision}
A submission commits the result and its evidence. A verifier checks the exact
agreement, allocation, input, output, checker and custody record before issuing
a decision. Missing evidence produces no decision. An available result that
fails the predicate produces rejection. These outcomes must remain distinct:
a verifier outage is not proof that the provider delivered bad work.

Off-chain artifacts use canonical JSON and Ed25519 signatures~\cite{jcs}.
The reference escrow accepts a separate EIP-712 decision~\cite{eip712}, binding
the allocation, agreement, submission commitment, decision digest, acceptance
bit, beneficiary, asset and amount. Its domain includes chain ID and contract
address. The contract verifies the pinned signer's decision, not the
application predicate. The publication artifact specifies which digests cover
bodies and which cover complete signed envelopes; those encodings cannot be
interchanged.

\subsection{A child is a separate obligation}
A parent link records provenance. It is not collateral and grants no authority.
Before hiring a specialist, the intermediary must commit funds it already
controls to a different allocation. In the opening example, the buyer deposits
100 and the intermediary deposits 60. The buyer's reserved 100 is not available
for the intermediary to pledge again. If the intermediary cannot finance the
60, this profile cannot admit the subcontract. Credit or conditional netting
would require a different contract and a different failure argument.

\begin{figure}[t]
\centering
\begin{tikzpicture}[font=\small,box/.style={draw=ink,rounded corners=2pt,align=center,minimum width=2.8cm,minimum height=0.9cm},>=Latex]
\node[box] (b) {Buyer $B$\\local authority};
\node[box,right=2.5cm of b] (p) {Provider $P$\\local authority};
\node[box,right=2.5cm of p] (s) {Specialist $S$\\local authority};
\draw[->,thick] (b) -- node[above]{agreement} (p);
\draw[->,thick] (p) -- node[above]{subcontract} (s);
\node[box,below=1.1cm of b] (f1) {Allocation $f_1$\\100 from $B$};
\node[box,below=1.1cm of s] (f2) {Allocation $f_2$\\60 from $P$};
\draw[->] (b) -- (f1);
\draw[->] (p.south) -- node[left]{deposit} (f2.north west);
\draw[->,dashed] (f2) -- node[right]{earned claim} (s);
\node[align=center,below=0.22cm of f1.south east,anchor=north west] {Parent refund leaves $f_2$ intact.\\Evidence crosses; each receiver admits locally.};
\end{tikzpicture}
\caption{Work dependencies and funding dependencies differ. Each obligation has
its own backing; the figure depicts neither a shared application operator nor
an absence of shared settlement and verification dependencies.}
\label{fig:work}
\end{figure}
